\documentclass[aspectratio=169,10pt]{beamer}
% ---------------------------------------------------------------
% R3.07 SQL dans un langage de programmation -- preambule commun des diapositives
% Charte derivee de R5.A.08 Qualite de developpement (metropolis, dark tech)
% Palette IUT partagee (violet du logo + accent or)
% Utilisation :
%   \documentclass[aspectratio=169,10pt]{beamer}
%   % ---------------------------------------------------------------
% R3.07 SQL dans un langage de programmation -- preambule commun des diapositives
% Charte derivee de R5.A.08 Qualite de developpement (metropolis, dark tech)
% Palette IUT partagee (violet du logo + accent or)
% Utilisation :
%   \documentclass[aspectratio=169,10pt]{beamer}
%   % ---------------------------------------------------------------
% R3.07 SQL dans un langage de programmation -- preambule commun des diapositives
% Charte derivee de R5.A.08 Qualite de developpement (metropolis, dark tech)
% Palette IUT partagee (violet du logo + accent or)
% Utilisation :
%   \documentclass[aspectratio=169,10pt]{beamer}
%   \input{../../latex/common/r307-slides-preamble.tex}
%   \graphicspath{{../../latex/common/assets/}{img/}{figs/}}
% ---------------------------------------------------------------
\usetheme{metropolis}
\usepackage{appendixnumberbeamer}

\usepackage[T1]{fontenc}
\usepackage[utf8]{inputenc}
\usepackage[french]{babel}
\usepackage[scale=0.94]{tgheros}
\renewcommand{\familydefault}{\sfdefault}
\usepackage{amsmath,amssymb,amsfonts}
\usepackage{graphicx}
\usepackage{booktabs}
\usepackage{multirow}
\usepackage{tikz}
\usetikzlibrary{shapes.geometric, shapes, arrows, arrows.meta, positioning, calc, fit}
\usepackage{xcolor}
\usepackage{listings}
\usepackage[skins,listings]{tcolorbox}
\usepackage{qrcode}
\usepackage{hyperref}

% ---------------------------------------------------------------
% PALETTE (dark-mode tech: near-black canvas, violet (logo IUT) + gold accents)
% ---------------------------------------------------------------
\definecolor{qabg}{HTML}{0E1016}      % canvas background
\definecolor{qapanel}{HTML}{181B24}   % card / block panel
\definecolor{qapanelalt}{HTML}{20242F}% secondary panel tint
\definecolor{qaline}{HTML}{2B303D}    % hairline rules / table rules
\definecolor{qatext}{HTML}{ECEEF2}    % primary text (off-white)
\definecolor{qamuted}{HTML}{8B93A6}   % secondary / caption text
\definecolor{qaaccent}{HTML}{D93ED7}  % primary accent -- violet du logo IUT, eclairci pour fond sombre
\definecolor{qaaccent2}{HTML}{F2B84B} % secondary accent -- or/ambre
\definecolor{qacodekw}{HTML}{C792EA}  % code : mots-cles
\definecolor{qacodestr}{HTML}{C3E88D} % code : chaines
\definecolor{qacodecom}{HTML}{7F8AA3} % code : commentaires
\definecolor{qacodeann}{HTML}{F2B84B} % code : annotations / javadoc

\setbeamercolor{background canvas}{bg=qabg}
\setbeamercolor{normal text}{fg=qatext}
\setbeamercolor{palette primary}{bg=qabg, fg=qamuted}
\setbeamercolor{frametitle}{bg=qabg, fg=qatext}
\setbeamercolor{progress bar}{fg=qaaccent, bg=qaline}
\setbeamercolor{alerted text}{fg=qaaccent2}
\setbeamercolor{block title}{bg=qapanel, fg=qaaccent}
\setbeamercolor{block body}{bg=qapanel, fg=qatext}
\setbeamercolor{block title alerted}{bg=qapanel, fg=qaaccent2}
\setbeamercolor{block body alerted}{bg=qapanel, fg=qatext}
\setbeamercolor{block title example}{bg=qapanelalt, fg=qaaccent}
\setbeamercolor{block body example}{bg=qapanelalt, fg=qatext}
\setbeamercolor{example text}{fg=qaaccent}
\setbeamercolor{itemize item}{fg=qaaccent}
\setbeamercolor{itemize subitem}{fg=qaaccent2}
\setbeamercolor{enumerate item}{fg=qaaccent}
\setbeamercolor{enumerate subitem}{fg=qaaccent2}
\setbeamercolor{caption}{fg=qamuted}
\setbeamercolor{page number in head/foot}{fg=qamuted, bg=qabg}
\setbeamercolor{title}{fg=qatext}
\setbeamercolor{subtitle}{fg=qamuted}
\setbeamercolor{author}{fg=qatext}
\setbeamercolor{institute}{fg=qamuted}
\setbeamercolor{date}{fg=qamuted}
\setbeamercolor{section in toc}{fg=qatext}
\setbeamercolor{subsection in toc}{fg=qamuted}

\setbeamertemplate{blocks}[rounded][shadow=false]
\setbeamertemplate{itemize item}{\raisebox{0.15em}{\scriptsize\textcolor{qaaccent}{\ensuremath{\blacktriangleright}}}}
\setbeamertemplate{itemize subitem}{\raisebox{0.1em}{\tiny\textcolor{qaaccent2}{\ensuremath{\bullet}}}}

\setbeamerfont{frametitle}{series=\bfseries, size=\Large}
\setbeamerfont{title}{series=\bfseries, size=\Large}
\setbeamerfont{subtitle}{size=\normalsize}
\setbeamerfont{institute}{size=\scriptsize}
\setbeamerfont{date}{size=\scriptsize}
\setbeamerfont{author}{size=\small}

\metroset{sectionpage=none, subsectionpage=none, numbering=fraction, progressbar=frametitle}
\setbeamertemplate{section in toc}[sections numbered]

% ---------------------------------------------------------------
% LISTINGS -- theme sombre (lstlisting reste utilisable tel quel)
% \lstcode : police + couleur de base, a reutiliser dans les
%            surcharges locales (ex. basicstyle=\lstcode\scriptsize)
% ---------------------------------------------------------------
\newcommand{\lstcode}{\ttfamily\color{qatext}}

\lstdefinelanguage{TypeScript}{
  keywords={typeof, new, true, false, catch, function, return, null, switch, var,
    if, in, while, do, else, case, break, const, let, interface, class, extends,
    implements, enum, export, import, this, super, constructor, public, private,
    protected, static, async, await, throw, throws, try, finally, void},
  ndkeywords={boolean, any, number, string, int, long, double},
  sensitive=true,
  comment=[l]{//},
  morecomment=[s]{/*}{*/},
  morestring=[b]',
  morestring=[b]"
}

\lstset{
  language=Java,
  basicstyle=\lstcode\footnotesize,
  keywordstyle=\color{qacodekw}\bfseries,
  ndkeywordstyle=\color{qaaccent2}\bfseries,
  identifierstyle=\color{qatext},
  stringstyle=\color{qacodestr},
  commentstyle=\color{qacodecom}\itshape,
  morecomment=[s][\color{qacodeann}]{/**}{*/},
  backgroundcolor=\color{qapanelalt},
  rulecolor=\color{qaline},
  frame=single,
  framesep=3pt,
  xleftmargin=4pt,
  xrightmargin=4pt,
  keepspaces=true,
  columns=fullflexible,
  showstringspaces=false,
  breaklines=true,
  tabsize=2,
  escapeinside={(*@}{@*)},
}

% ---------------------------------------------------------------
% CARD-STYLE BOXES
% ---------------------------------------------------------------
\tcbset{
  figcard/.style={colback=white, colframe=white, boxrule=0pt, arc=4pt,
    left=3pt, right=3pt, top=3pt, bottom=3pt, boxsep=0pt,
    halign=center, valign=center},
  datacard/.style={colback=qapanel, colframe=qaline, boxrule=0.4pt, arc=3pt,
    left=8pt, right=8pt, top=6pt, bottom=6pt, coltext=qatext,
    before upper={\color{qatext}}, halign=center},
  codebox/.style={colback=qapanelalt, colframe=qaline, boxrule=0.4pt, arc=3pt,
    left=8pt, right=8pt, top=5pt, bottom=5pt, listing only,
    listing options={basicstyle=\ttfamily\small\color{qaaccent},
      breaklines=true, showstringspaces=false, columns=fullflexible,
      keepspaces=true, frame=none, backgroundcolor={}}}
}
\newtcolorbox{figcard}[1][]{figcard,#1}
\newtcolorbox{datacard}{datacard}
\newtcblisting{codebox}{codebox}

% ---------------------------------------------------------------
% HANDY MACROS
% ---------------------------------------------------------------
\newcommand{\hl}[1]{\alert{#1}}
\newcommand{\source}[1]{{\scriptsize\color{qamuted}#1}}

% Sommaire automatique en debut de section
\AtBeginSection[]{%
  \begin{frame}<beamer>{Plan du cours}
    \tableofcontents[currentsection]
  \end{frame}}

% Page de titre : \rtitleslide{numero de chapitre}{titre du chapitre}
% Les auteurs sont repris de \rauthors (modifiable avant \begin{document})
\newcommand{\rauthors}{Selena Lamari et Chouki Tibermacine}
\newcommand{\rtitleslide}[2]{%
\begin{frame}[plain,noframenumbering]
  \begin{center}
    \vspace{2mm}
    {\color{qaaccent}\rule{2.1cm}{1.2pt}}\\[6mm]
    {\Large\bfseries\color{qatext} R3.07}\\[1mm]
    {\Large\bfseries\color{qatext} SQL dans un langage de programmation}\\[5mm]
    {\large\color{qamuted} Chapitre #1}\\[1mm]
    {\large\color{qaaccent} #2}\\[6mm]
    \includegraphics[height=2.1cm]{iut_logo_transparent.png}

    \vspace{4mm}

    {\small\color{qatext} \rauthors}\\[1mm]
    {\scriptsize\color{qamuted} BUT 2 Informatique --- IUT de Vannes --- Universit\'e Bretagne Sud}
  \end{center}
\end{frame}
}

% ---- Diapo de fin « Questions ? » (même modèle que R5.A.08) ----
\newcommand{\rcontact}{Chouki.Tibermacine@univ-ubs.fr}
\newcommand{\rquestionslide}{%
\begin{frame}{Questions ?}
  \begin{center}
    \vspace{2mm}
    \begin{tikzpicture}
      \node[regular polygon, regular polygon sides=6, draw=qaline, thick,
        minimum size=4.6cm, rotate=90] (hexin) {};
      \node[regular polygon, regular polygon sides=6, draw=qaaccent, dashed,
        thick, minimum size=6.2cm, rotate=90] (hexout) {};
      \node[font=\fontsize{92}{92}\selectfont\bfseries, color=qaaccent2] at (hexin) {?};
    \end{tikzpicture}\\[5mm]
    {\small\color{qamuted} \rcontact}
  \end{center}
\end{frame}}

%   \graphicspath{{../../latex/common/assets/}{img/}{figs/}}
% ---------------------------------------------------------------
\usetheme{metropolis}
\usepackage{appendixnumberbeamer}

\usepackage[T1]{fontenc}
\usepackage[utf8]{inputenc}
\usepackage[french]{babel}
\usepackage[scale=0.94]{tgheros}
\renewcommand{\familydefault}{\sfdefault}
\usepackage{amsmath,amssymb,amsfonts}
\usepackage{graphicx}
\usepackage{booktabs}
\usepackage{multirow}
\usepackage{tikz}
\usetikzlibrary{shapes.geometric, shapes, arrows, arrows.meta, positioning, calc, fit}
\usepackage{xcolor}
\usepackage{listings}
\usepackage[skins,listings]{tcolorbox}
\usepackage{qrcode}
\usepackage{hyperref}

% ---------------------------------------------------------------
% PALETTE (dark-mode tech: near-black canvas, violet (logo IUT) + gold accents)
% ---------------------------------------------------------------
\definecolor{qabg}{HTML}{0E1016}      % canvas background
\definecolor{qapanel}{HTML}{181B24}   % card / block panel
\definecolor{qapanelalt}{HTML}{20242F}% secondary panel tint
\definecolor{qaline}{HTML}{2B303D}    % hairline rules / table rules
\definecolor{qatext}{HTML}{ECEEF2}    % primary text (off-white)
\definecolor{qamuted}{HTML}{8B93A6}   % secondary / caption text
\definecolor{qaaccent}{HTML}{D93ED7}  % primary accent -- violet du logo IUT, eclairci pour fond sombre
\definecolor{qaaccent2}{HTML}{F2B84B} % secondary accent -- or/ambre
\definecolor{qacodekw}{HTML}{C792EA}  % code : mots-cles
\definecolor{qacodestr}{HTML}{C3E88D} % code : chaines
\definecolor{qacodecom}{HTML}{7F8AA3} % code : commentaires
\definecolor{qacodeann}{HTML}{F2B84B} % code : annotations / javadoc

\setbeamercolor{background canvas}{bg=qabg}
\setbeamercolor{normal text}{fg=qatext}
\setbeamercolor{palette primary}{bg=qabg, fg=qamuted}
\setbeamercolor{frametitle}{bg=qabg, fg=qatext}
\setbeamercolor{progress bar}{fg=qaaccent, bg=qaline}
\setbeamercolor{alerted text}{fg=qaaccent2}
\setbeamercolor{block title}{bg=qapanel, fg=qaaccent}
\setbeamercolor{block body}{bg=qapanel, fg=qatext}
\setbeamercolor{block title alerted}{bg=qapanel, fg=qaaccent2}
\setbeamercolor{block body alerted}{bg=qapanel, fg=qatext}
\setbeamercolor{block title example}{bg=qapanelalt, fg=qaaccent}
\setbeamercolor{block body example}{bg=qapanelalt, fg=qatext}
\setbeamercolor{example text}{fg=qaaccent}
\setbeamercolor{itemize item}{fg=qaaccent}
\setbeamercolor{itemize subitem}{fg=qaaccent2}
\setbeamercolor{enumerate item}{fg=qaaccent}
\setbeamercolor{enumerate subitem}{fg=qaaccent2}
\setbeamercolor{caption}{fg=qamuted}
\setbeamercolor{page number in head/foot}{fg=qamuted, bg=qabg}
\setbeamercolor{title}{fg=qatext}
\setbeamercolor{subtitle}{fg=qamuted}
\setbeamercolor{author}{fg=qatext}
\setbeamercolor{institute}{fg=qamuted}
\setbeamercolor{date}{fg=qamuted}
\setbeamercolor{section in toc}{fg=qatext}
\setbeamercolor{subsection in toc}{fg=qamuted}

\setbeamertemplate{blocks}[rounded][shadow=false]
\setbeamertemplate{itemize item}{\raisebox{0.15em}{\scriptsize\textcolor{qaaccent}{\ensuremath{\blacktriangleright}}}}
\setbeamertemplate{itemize subitem}{\raisebox{0.1em}{\tiny\textcolor{qaaccent2}{\ensuremath{\bullet}}}}

\setbeamerfont{frametitle}{series=\bfseries, size=\Large}
\setbeamerfont{title}{series=\bfseries, size=\Large}
\setbeamerfont{subtitle}{size=\normalsize}
\setbeamerfont{institute}{size=\scriptsize}
\setbeamerfont{date}{size=\scriptsize}
\setbeamerfont{author}{size=\small}

\metroset{sectionpage=none, subsectionpage=none, numbering=fraction, progressbar=frametitle}
\setbeamertemplate{section in toc}[sections numbered]

% ---------------------------------------------------------------
% LISTINGS -- theme sombre (lstlisting reste utilisable tel quel)
% \lstcode : police + couleur de base, a reutiliser dans les
%            surcharges locales (ex. basicstyle=\lstcode\scriptsize)
% ---------------------------------------------------------------
\newcommand{\lstcode}{\ttfamily\color{qatext}}

\lstdefinelanguage{TypeScript}{
  keywords={typeof, new, true, false, catch, function, return, null, switch, var,
    if, in, while, do, else, case, break, const, let, interface, class, extends,
    implements, enum, export, import, this, super, constructor, public, private,
    protected, static, async, await, throw, throws, try, finally, void},
  ndkeywords={boolean, any, number, string, int, long, double},
  sensitive=true,
  comment=[l]{//},
  morecomment=[s]{/*}{*/},
  morestring=[b]',
  morestring=[b]"
}

\lstset{
  language=Java,
  basicstyle=\lstcode\footnotesize,
  keywordstyle=\color{qacodekw}\bfseries,
  ndkeywordstyle=\color{qaaccent2}\bfseries,
  identifierstyle=\color{qatext},
  stringstyle=\color{qacodestr},
  commentstyle=\color{qacodecom}\itshape,
  morecomment=[s][\color{qacodeann}]{/**}{*/},
  backgroundcolor=\color{qapanelalt},
  rulecolor=\color{qaline},
  frame=single,
  framesep=3pt,
  xleftmargin=4pt,
  xrightmargin=4pt,
  keepspaces=true,
  columns=fullflexible,
  showstringspaces=false,
  breaklines=true,
  tabsize=2,
  escapeinside={(*@}{@*)},
}

% ---------------------------------------------------------------
% CARD-STYLE BOXES
% ---------------------------------------------------------------
\tcbset{
  figcard/.style={colback=white, colframe=white, boxrule=0pt, arc=4pt,
    left=3pt, right=3pt, top=3pt, bottom=3pt, boxsep=0pt,
    halign=center, valign=center},
  datacard/.style={colback=qapanel, colframe=qaline, boxrule=0.4pt, arc=3pt,
    left=8pt, right=8pt, top=6pt, bottom=6pt, coltext=qatext,
    before upper={\color{qatext}}, halign=center},
  codebox/.style={colback=qapanelalt, colframe=qaline, boxrule=0.4pt, arc=3pt,
    left=8pt, right=8pt, top=5pt, bottom=5pt, listing only,
    listing options={basicstyle=\ttfamily\small\color{qaaccent},
      breaklines=true, showstringspaces=false, columns=fullflexible,
      keepspaces=true, frame=none, backgroundcolor={}}}
}
\newtcolorbox{figcard}[1][]{figcard,#1}
\newtcolorbox{datacard}{datacard}
\newtcblisting{codebox}{codebox}

% ---------------------------------------------------------------
% HANDY MACROS
% ---------------------------------------------------------------
\newcommand{\hl}[1]{\alert{#1}}
\newcommand{\source}[1]{{\scriptsize\color{qamuted}#1}}

% Sommaire automatique en debut de section
\AtBeginSection[]{%
  \begin{frame}<beamer>{Plan du cours}
    \tableofcontents[currentsection]
  \end{frame}}

% Page de titre : \rtitleslide{numero de chapitre}{titre du chapitre}
% Les auteurs sont repris de \rauthors (modifiable avant \begin{document})
\newcommand{\rauthors}{Selena Lamari et Chouki Tibermacine}
\newcommand{\rtitleslide}[2]{%
\begin{frame}[plain,noframenumbering]
  \begin{center}
    \vspace{2mm}
    {\color{qaaccent}\rule{2.1cm}{1.2pt}}\\[6mm]
    {\Large\bfseries\color{qatext} R3.07}\\[1mm]
    {\Large\bfseries\color{qatext} SQL dans un langage de programmation}\\[5mm]
    {\large\color{qamuted} Chapitre #1}\\[1mm]
    {\large\color{qaaccent} #2}\\[6mm]
    \includegraphics[height=2.1cm]{iut_logo_transparent.png}

    \vspace{4mm}

    {\small\color{qatext} \rauthors}\\[1mm]
    {\scriptsize\color{qamuted} BUT 2 Informatique --- IUT de Vannes --- Universit\'e Bretagne Sud}
  \end{center}
\end{frame}
}

% ---- Diapo de fin « Questions ? » (même modèle que R5.A.08) ----
\newcommand{\rcontact}{Chouki.Tibermacine@univ-ubs.fr}
\newcommand{\rquestionslide}{%
\begin{frame}{Questions ?}
  \begin{center}
    \vspace{2mm}
    \begin{tikzpicture}
      \node[regular polygon, regular polygon sides=6, draw=qaline, thick,
        minimum size=4.6cm, rotate=90] (hexin) {};
      \node[regular polygon, regular polygon sides=6, draw=qaaccent, dashed,
        thick, minimum size=6.2cm, rotate=90] (hexout) {};
      \node[font=\fontsize{92}{92}\selectfont\bfseries, color=qaaccent2] at (hexin) {?};
    \end{tikzpicture}\\[5mm]
    {\small\color{qamuted} \rcontact}
  \end{center}
\end{frame}}

%   \graphicspath{{../../latex/common/assets/}{img/}{figs/}}
% ---------------------------------------------------------------
\usetheme{metropolis}
\usepackage{appendixnumberbeamer}

\usepackage[T1]{fontenc}
\usepackage[utf8]{inputenc}
\usepackage[french]{babel}
\usepackage[scale=0.94]{tgheros}
\renewcommand{\familydefault}{\sfdefault}
\usepackage{amsmath,amssymb,amsfonts}
\usepackage{graphicx}
\usepackage{booktabs}
\usepackage{multirow}
\usepackage{tikz}
\usetikzlibrary{shapes.geometric, shapes, arrows, arrows.meta, positioning, calc, fit}
\usepackage{xcolor}
\usepackage{listings}
\usepackage[skins,listings]{tcolorbox}
\usepackage{qrcode}
\usepackage{hyperref}

% ---------------------------------------------------------------
% PALETTE (dark-mode tech: near-black canvas, violet (logo IUT) + gold accents)
% ---------------------------------------------------------------
\definecolor{qabg}{HTML}{0E1016}      % canvas background
\definecolor{qapanel}{HTML}{181B24}   % card / block panel
\definecolor{qapanelalt}{HTML}{20242F}% secondary panel tint
\definecolor{qaline}{HTML}{2B303D}    % hairline rules / table rules
\definecolor{qatext}{HTML}{ECEEF2}    % primary text (off-white)
\definecolor{qamuted}{HTML}{8B93A6}   % secondary / caption text
\definecolor{qaaccent}{HTML}{D93ED7}  % primary accent -- violet du logo IUT, eclairci pour fond sombre
\definecolor{qaaccent2}{HTML}{F2B84B} % secondary accent -- or/ambre
\definecolor{qacodekw}{HTML}{C792EA}  % code : mots-cles
\definecolor{qacodestr}{HTML}{C3E88D} % code : chaines
\definecolor{qacodecom}{HTML}{7F8AA3} % code : commentaires
\definecolor{qacodeann}{HTML}{F2B84B} % code : annotations / javadoc

\setbeamercolor{background canvas}{bg=qabg}
\setbeamercolor{normal text}{fg=qatext}
\setbeamercolor{palette primary}{bg=qabg, fg=qamuted}
\setbeamercolor{frametitle}{bg=qabg, fg=qatext}
\setbeamercolor{progress bar}{fg=qaaccent, bg=qaline}
\setbeamercolor{alerted text}{fg=qaaccent2}
\setbeamercolor{block title}{bg=qapanel, fg=qaaccent}
\setbeamercolor{block body}{bg=qapanel, fg=qatext}
\setbeamercolor{block title alerted}{bg=qapanel, fg=qaaccent2}
\setbeamercolor{block body alerted}{bg=qapanel, fg=qatext}
\setbeamercolor{block title example}{bg=qapanelalt, fg=qaaccent}
\setbeamercolor{block body example}{bg=qapanelalt, fg=qatext}
\setbeamercolor{example text}{fg=qaaccent}
\setbeamercolor{itemize item}{fg=qaaccent}
\setbeamercolor{itemize subitem}{fg=qaaccent2}
\setbeamercolor{enumerate item}{fg=qaaccent}
\setbeamercolor{enumerate subitem}{fg=qaaccent2}
\setbeamercolor{caption}{fg=qamuted}
\setbeamercolor{page number in head/foot}{fg=qamuted, bg=qabg}
\setbeamercolor{title}{fg=qatext}
\setbeamercolor{subtitle}{fg=qamuted}
\setbeamercolor{author}{fg=qatext}
\setbeamercolor{institute}{fg=qamuted}
\setbeamercolor{date}{fg=qamuted}
\setbeamercolor{section in toc}{fg=qatext}
\setbeamercolor{subsection in toc}{fg=qamuted}

\setbeamertemplate{blocks}[rounded][shadow=false]
\setbeamertemplate{itemize item}{\raisebox{0.15em}{\scriptsize\textcolor{qaaccent}{\ensuremath{\blacktriangleright}}}}
\setbeamertemplate{itemize subitem}{\raisebox{0.1em}{\tiny\textcolor{qaaccent2}{\ensuremath{\bullet}}}}

\setbeamerfont{frametitle}{series=\bfseries, size=\Large}
\setbeamerfont{title}{series=\bfseries, size=\Large}
\setbeamerfont{subtitle}{size=\normalsize}
\setbeamerfont{institute}{size=\scriptsize}
\setbeamerfont{date}{size=\scriptsize}
\setbeamerfont{author}{size=\small}

\metroset{sectionpage=none, subsectionpage=none, numbering=fraction, progressbar=frametitle}
\setbeamertemplate{section in toc}[sections numbered]

% ---------------------------------------------------------------
% LISTINGS -- theme sombre (lstlisting reste utilisable tel quel)
% \lstcode : police + couleur de base, a reutiliser dans les
%            surcharges locales (ex. basicstyle=\lstcode\scriptsize)
% ---------------------------------------------------------------
\newcommand{\lstcode}{\ttfamily\color{qatext}}

\lstdefinelanguage{TypeScript}{
  keywords={typeof, new, true, false, catch, function, return, null, switch, var,
    if, in, while, do, else, case, break, const, let, interface, class, extends,
    implements, enum, export, import, this, super, constructor, public, private,
    protected, static, async, await, throw, throws, try, finally, void},
  ndkeywords={boolean, any, number, string, int, long, double},
  sensitive=true,
  comment=[l]{//},
  morecomment=[s]{/*}{*/},
  morestring=[b]',
  morestring=[b]"
}

\lstset{
  language=Java,
  basicstyle=\lstcode\footnotesize,
  keywordstyle=\color{qacodekw}\bfseries,
  ndkeywordstyle=\color{qaaccent2}\bfseries,
  identifierstyle=\color{qatext},
  stringstyle=\color{qacodestr},
  commentstyle=\color{qacodecom}\itshape,
  morecomment=[s][\color{qacodeann}]{/**}{*/},
  backgroundcolor=\color{qapanelalt},
  rulecolor=\color{qaline},
  frame=single,
  framesep=3pt,
  xleftmargin=4pt,
  xrightmargin=4pt,
  keepspaces=true,
  columns=fullflexible,
  showstringspaces=false,
  breaklines=true,
  tabsize=2,
  escapeinside={(*@}{@*)},
}

% ---------------------------------------------------------------
% CARD-STYLE BOXES
% ---------------------------------------------------------------
\tcbset{
  figcard/.style={colback=white, colframe=white, boxrule=0pt, arc=4pt,
    left=3pt, right=3pt, top=3pt, bottom=3pt, boxsep=0pt,
    halign=center, valign=center},
  datacard/.style={colback=qapanel, colframe=qaline, boxrule=0.4pt, arc=3pt,
    left=8pt, right=8pt, top=6pt, bottom=6pt, coltext=qatext,
    before upper={\color{qatext}}, halign=center},
  codebox/.style={colback=qapanelalt, colframe=qaline, boxrule=0.4pt, arc=3pt,
    left=8pt, right=8pt, top=5pt, bottom=5pt, listing only,
    listing options={basicstyle=\ttfamily\small\color{qaaccent},
      breaklines=true, showstringspaces=false, columns=fullflexible,
      keepspaces=true, frame=none, backgroundcolor={}}}
}
\newtcolorbox{figcard}[1][]{figcard,#1}
\newtcolorbox{datacard}{datacard}
\newtcblisting{codebox}{codebox}

% ---------------------------------------------------------------
% HANDY MACROS
% ---------------------------------------------------------------
\newcommand{\hl}[1]{\alert{#1}}
\newcommand{\source}[1]{{\scriptsize\color{qamuted}#1}}

% Sommaire automatique en debut de section
\AtBeginSection[]{%
  \begin{frame}<beamer>{Plan du cours}
    \tableofcontents[currentsection]
  \end{frame}}

% Page de titre : \rtitleslide{numero de chapitre}{titre du chapitre}
% Les auteurs sont repris de \rauthors (modifiable avant \begin{document})
\newcommand{\rauthors}{Selena Lamari et Chouki Tibermacine}
\newcommand{\rtitleslide}[2]{%
\begin{frame}[plain,noframenumbering]
  \begin{center}
    \vspace{2mm}
    {\color{qaaccent}\rule{2.1cm}{1.2pt}}\\[6mm]
    {\Large\bfseries\color{qatext} R3.07}\\[1mm]
    {\Large\bfseries\color{qatext} SQL dans un langage de programmation}\\[5mm]
    {\large\color{qamuted} Chapitre #1}\\[1mm]
    {\large\color{qaaccent} #2}\\[6mm]
    \includegraphics[height=2.1cm]{iut_logo_transparent.png}

    \vspace{4mm}

    {\small\color{qatext} \rauthors}\\[1mm]
    {\scriptsize\color{qamuted} BUT 2 Informatique --- IUT de Vannes --- Universit\'e Bretagne Sud}
  \end{center}
\end{frame}
}

% ---- Diapo de fin « Questions ? » (même modèle que R5.A.08) ----
\newcommand{\rcontact}{Chouki.Tibermacine@univ-ubs.fr}
\newcommand{\rquestionslide}{%
\begin{frame}{Questions ?}
  \begin{center}
    \vspace{2mm}
    \begin{tikzpicture}
      \node[regular polygon, regular polygon sides=6, draw=qaline, thick,
        minimum size=4.6cm, rotate=90] (hexin) {};
      \node[regular polygon, regular polygon sides=6, draw=qaaccent, dashed,
        thick, minimum size=6.2cm, rotate=90] (hexout) {};
      \node[font=\fontsize{92}{92}\selectfont\bfseries, color=qaaccent2] at (hexin) {?};
    \end{tikzpicture}\\[5mm]
    {\small\color{qamuted} \rcontact}
  \end{center}
\end{frame}}

\graphicspath{{../../latex/common/assets/}}

\title{Index et optimisation}
\author[C. Tibermacine]{Selena Lamari et Chouki Tibermacine}
\date{}

\lstdefinestyle{plan}{language={}, basicstyle=\lstcode\scriptsize, keywordstyle={}, commentstyle={}, stringstyle={}}

\tikzset{
  bnode/.style={rectangle, rounded corners=2pt, draw=qaaccent, thick, fill=qapanel, text=qatext,
    minimum height=0.5cm, align=center, font=\scriptsize\ttfamily},
  leaf/.style={bnode, draw=qaaccent2},
  row/.style={rectangle, draw=qaline, fill=qapanelalt, text=qamuted, minimum width=0.9cm,
    minimum height=0.4cm, font=\tiny\ttfamily}
}

\begin{document}

\rtitleslide{3}{Index et optimisation}

% ================================================================
\section{Comment le SGBD ex\'ecute une requ\^ete}

\begin{frame}{De la requ\^ete SQL au r\'esultat}
  \begin{itemize}
    \item SQL est \textbf{d\'eclaratif} : on dit \emph{quoi}, pas \emph{comment}
    \item L'\textbf{optimiseur} (\textit{planner}) choisit un \textbf{plan d'ex\'ecution} parmi de nombreux possibles :
      \begin{itemize}
        \item parcourir toute la table (\textit{Seq Scan}) ou passer par un index (\textit{Index Scan}, \textit{Bitmap Scan}) ?
        \item quel algorithme de jointure (\textit{Nested Loop}, \textit{Hash Join}, \textit{Merge Join}) ? dans quel ordre ?
      \end{itemize}
    \item Il estime un \textbf{co\^ut} pour chaque plan \`a partir de \textbf{statistiques} sur les donn\'ees (\texttt{ANALYZE})
    \item Le d\'eveloppeur agit sur deux leviers :
      \begin{enumerate}
        \item la \textbf{structure physique} : index, types, contraintes
        \item la \textbf{fa\c{c}on d'interroger} depuis l'application : requ\^etes, lots, pagination, nombre d'allers-retours
      \end{enumerate}
  \end{itemize}
\end{frame}

\begin{frame}[fragile]{Lire un plan : \texttt{EXPLAIN ANALYZE}}
  Table \texttt{Operation} : 1 million de lignes, aucun index sur \texttt{numCompte}
\begin{lstlisting}[style=plan]
EXPLAIN ANALYZE SELECT * FROM Operation WHERE numCompte = 4242;

 Gather  (cost=1000.00..13561.93 rows=6 width=25) (actual time=5.211..18.546 rows=2 loops=1)
   Workers Planned: 2
   Workers Launched: 2
   ->  Parallel Seq Scan on operation  (cost=0.00..12561.33 rows=2 width=25)
                                       (actual time=10.874..14.583 rows=1 loops=3)
         Filter: (numcompte = 4242)
         Rows Removed by Filter: 333333
 Planning Time: 0.284 ms
 Execution Time: 18.609 ms
\end{lstlisting}
  \begin{itemize}
    \item \texttt{EXPLAIN} : plan \textbf{estim\'e} ; \texttt{EXPLAIN ANALYZE} : ex\'ecute r\'eellement et mesure
    \item \texttt{cost} : unit\'e arbitraire ; \texttt{rows} estim\'ees vs r\'eelles ; \texttt{Rows Removed by Filter} : travail gaspill\'e
    \item Ici : \alert{1 million de lignes lues pour en garder 2}
  \end{itemize}
\end{frame}

% ================================================================
\section{Les index}

\begin{frame}{Un index B-tree}
  \begin{columns}[c]
    \begin{column}{0.55\textwidth}
      \begin{center}
      \begin{tikzpicture}[level distance=1.0cm, sibling distance=2.1cm]
        \node[bnode] {2000 | 4000}
          child[qamuted] {node[leaf] {500 | 1200}}
          child[qamuted] {node[leaf] {2500 | 3100}}
          child[qamuted] {node[leaf] {4242 | 4800}};
        \node[row] (r1) at (-1.8,-2.2) {ligne};
        \node[row] (r2) at (0,-2.2) {ligne};
        \node[row] (r3) at (1.8,-2.2) {ligne};
        \draw[-Latex, qaaccent2, thick] (2.1,-1.25) -- (r3.north);
        \node[font=\tiny, text=qamuted] at (0,-2.75) {table (tas de lignes, non tri\'ees)};
      \end{tikzpicture}
      \end{center}
    \end{column}
    \begin{column}{0.42\textwidth}
      \begin{itemize}
        \item Arbre \textbf{tri\'e et \'equilibr\'e} sur la valeur de la colonne
        \item Recherche en $O(\log n)$ au lieu de $O(n)$
        \item Sert pour \texttt{=}, \texttt{<}, \texttt{>}, \texttt{BETWEEN}, \texttt{ORDER BY}, jointures
        \item Type par d\'efaut de \texttt{CREATE INDEX}
      \end{itemize}
    \end{column}
  \end{columns}
\end{frame}

\begin{frame}[fragile]{Cr\'eer un index et mesurer}
\begin{lstlisting}[language=SQL]
CREATE INDEX idx_operation_compte ON Operation(numCompte);
\end{lstlisting}
\begin{lstlisting}[style=plan]
 Bitmap Heap Scan on operation  (cost=4.47..28.03 rows=6 width=25) (actual time=0.017..0.019 rows=2 loops=1)
   Recheck Cond: (numcompte = 4242)
   ->  Bitmap Index Scan on idx_operation_compte  (cost=0.00..4.47 rows=6 width=0)
                                                  (actual time=0.013..0.013 rows=2 loops=1)
         Index Cond: (numcompte = 4242)
 Execution Time: 0.029 ms
\end{lstlisting}
  \begin{datacard}
    \small
    \begin{tabular}{@{}lrr@{}}
      \toprule
      \textbf{Requ\^ete (mesures PostgreSQL 16, 1 M d'op\'erations)} & \textbf{sans index} & \textbf{avec index} \\
      \midrule
      Op\'erations d'un compte & 19 ms & 0,03 ms \\
      Somme des op\'erations des comptes d'un client & 39 ms & 0,06 ms \\
      \texttt{DELETE} de 50 comptes (v\'erif. de cl\'e \'etrang\`ere) & \alert{2\,182 ms} & 0,6 ms \\
      \bottomrule
    \end{tabular}
  \end{datacard}
\end{frame}

\begin{frame}{Quels index existent d\'ej\`a ?}
  \begin{itemize}
    \item PostgreSQL cr\'ee \textbf{automatiquement} un index pour chaque \texttt{PRIMARY KEY} et chaque contrainte \texttt{UNIQUE}
    \item Il \alert{ne cr\'ee pas} d'index sur les \textbf{cl\'es \'etrang\`eres} :
      \begin{itemize}
        \item jointures et recherches \og les op\'erations du compte X \fg{} : parcours complet
        \item \texttt{DELETE}/\texttt{UPDATE} sur la table r\'ef\'erenc\'ee : pour chaque ligne supprim\'ee, la base v\'erifie la table fille\ldots{} en la parcourant enti\`erement
      \end{itemize}
    \item[$\Rightarrow$] R\`egle pratique : \textbf{indexer les colonnes de cl\'e \'etrang\`ere} utilis\'ees dans les jointures
    \item Lister les index : \texttt{\textbackslash di} dans \texttt{psql}, ou la vue \texttt{pg\_indexes}
  \end{itemize}
\end{frame}

\begin{frame}[fragile]{Quand l'index n'est pas utilis\'e}
  \begin{columns}[T]
    \begin{column}{0.5\textwidth}
      \textbf{Fonction appliqu\'ee \`a la colonne}
\begin{lstlisting}[language=SQL]
-- index sur nomClient ignore :
WHERE lower(nomClient) = 'dupont'
-- solution : index sur l'expression
CREATE INDEX ON Client(lower(nomClient));
\end{lstlisting}
      \textbf{\texttt{LIKE} par pr\'efixe}
\begin{lstlisting}[language=SQL]
-- ignore avec une collation fr/en :
WHERE nomClient LIKE 'Dup%'
CREATE INDEX ON Client
  (nomClient varchar_pattern_ops);
-- '%pont' : aucun B-tree ne peut aider
\end{lstlisting}
    \end{column}
    \begin{column}{0.47\textwidth}
      \textbf{Index composite : l'ordre compte}
\begin{lstlisting}[language=SQL]
CREATE INDEX ON Operation
  (numCompte, dateOperation);
-- utilise :
WHERE numCompte = 3 AND dateOperation > ?
WHERE numCompte = 3
-- PAS utilise (2e colonne seule) :
WHERE dateOperation = ?
\end{lstlisting}
      \textbf{Faible s\'electivit\'e}
      \begin{itemize}
        \item \texttt{typeOperation = 'DEPOT'} : 1 ligne sur 3
        \item l'index aide peu (54 ms $\to$ 28 ms) ; l'optimiseur pr\'ef\`ere parfois le parcours complet
      \end{itemize}
    \end{column}
  \end{columns}
\end{frame}

\begin{frame}{Un index a un co\^ut}
  \begin{itemize}
    \item Chaque \texttt{INSERT}, \texttt{UPDATE}, \texttt{DELETE} doit aussi \textbf{mettre \`a jour tous les index} de la table
    \item Mesure : insertion de 300\,000 op\'erations
      \begin{itemize}
        \item sans index secondaire : \textbf{0,13 s}
        \item avec 4 index secondaires : \alert{\textbf{1,3 s}} (10 fois plus lent)
      \end{itemize}
    \item Les index occupent de l'espace disque et de la m\'emoire
  \end{itemize}
  \begin{block}{D\'emarche}
    \begin{enumerate}
      \item identifier les requ\^etes \textbf{fr\'equentes} ou \textbf{lentes} de l'application
      \item les analyser avec \texttt{EXPLAIN ANALYZE}
      \item cr\'eer l'index qui correspond \textbf{exactement} \`a la clause \texttt{WHERE}/\texttt{JOIN}/\texttt{ORDER BY}
      \item re-mesurer ; supprimer les index inutiles
    \end{enumerate}
  \end{block}
\end{frame}

% ================================================================
\section{Optimiser c\^ot\'e application}

\begin{frame}[fragile]{R\'eduire les allers-retours : les lots (\textit{batch})}
\begin{lstlisting}[language=Java]
c.setAutoCommit(false);                       // une seule transaction
try (PreparedStatement ps = c.prepareStatement(
      "INSERT INTO Operation VALUES (?, ?, ?, ?, ?)")) {
  for (Operation op : operations) {
    ps.setInt(1, op.getNum()); /* ... */
    ps.addBatch();                             // empile au lieu d'envoyer
    if (++n % 1000 == 0) ps.executeBatch();    // envoie par paquets de 1000
  }
  ps.executeBatch();
}
c.commit();
\end{lstlisting}
  \begin{itemize}
    \item Une requ\^ete = un aller-retour r\'eseau + (en autocommit) une \'ecriture disque synchrone
    \item Avec le pilote PostgreSQL, ajouter \texttt{?reWriteBatchedInserts=true} \`a l'URL JDBC
      regroupe les lignes en \texttt{INSERT} multi-valeurs
    \item Gain typique : un \textbf{ordre de grandeur} ou plus (\`a mesurer en TP)
  \end{itemize}
\end{frame}

\begin{frame}[fragile]{Pagination : \texttt{OFFSET} vs pagination par cl\'e}
  \begin{columns}[T]
    \begin{column}{0.48\textwidth}
\begin{lstlisting}[language=SQL]
-- page 45 001 : lit et jette 900 000 lignes
SELECT * FROM Operation
 ORDER BY numOperation
 LIMIT 20 OFFSET 900000;
\end{lstlisting}
      \centering\alert{79 ms}
    \end{column}
    \begin{column}{0.48\textwidth}
\begin{lstlisting}[language=SQL]
-- on part de la derniere cle affichee
SELECT * FROM Operation
 WHERE numOperation > 900000
 ORDER BY numOperation LIMIT 20;
\end{lstlisting}
      \centering\textbf{0,02 ms}
    \end{column}
  \end{columns}
  \vspace{3mm}
  \begin{itemize}
    \item \texttt{OFFSET} est simple mais son co\^ut cro\^it avec le num\'ero de page
    \item La pagination par cl\'e (\textit{keyset}) utilise l'index et reste constante
  \end{itemize}
\end{frame}

\begin{frame}{Autres r\`egles d'hygi\`ene}
  \begin{itemize}
    \item \texttt{PreparedStatement} partout : s\'ecurit\'e (injection SQL) et r\'eutilisation du plan par le serveur
    \item S\'electionner \textbf{les colonnes utiles} plut\^ot que \texttt{SELECT *}
    \item Filtrer, trier, agr\'eger \textbf{en SQL} plut\^ot qu'en Java (ne pas rapatrier 1 M de lignes pour en compter 10)
    \item \'Eviter les requ\^etes dans une boucle : le probl\`eme \og \textbf{N+1} \fg{} (on le retrouvera avec JPA : TP4 et chapitre 6)
    \item R\'eutiliser les connexions : \textbf{pool} (HikariCP)
    \item Mesurer avant d'optimiser : \texttt{EXPLAIN ANALYZE}, \texttt{System.nanoTime()}, \texttt{pg\_stat\_statements}
  \end{itemize}
\end{frame}

\rquestionslide

\end{document}
